%---
% title: "Six synthetic curricula under two controls"
% date: 2026-07-10
% work_order: 5
% permalink: /posts/2026/07/prepretraining-tested-it-doesnt-work/
% tags:
%   - machine-learning
%   - language-models
%   - pretraining
%   - experiments
%---
\documentclass[11pt]{article}
\usepackage{publication-web}
\begin{document}

Prepretraining trains a language model on synthetic or abstract data before continuing on natural text. Its transfer claim is that the synthetic stage yields a better initialization than either random initialization or the same upstream training budget on natural text. A shared \href{https://github.com/Axym-Labs/pptrain}{\texttt{pptrain}} protocol tested six synthetic curricula---NCA sequences, LIME reasoning tasks, simpler symbolic transformations, procedural programs, Dyck languages, and synthetic summarization---on a Pythia-410M architecture with three seeds (\href{https://arxiv.org/abs/2603.10055}{Lee et al.~2026}; \href{https://proceedings.mlr.press/v139/wu21c.html}{Wu et al.~2021}; \href{https://arxiv.org/abs/2206.10139}{Wu et al.~2022}; \href{https://arxiv.org/abs/2601.21725}{Jiang et al.~2026}; \href{https://aclanthology.org/2021.findings-emnlp.273/}{Krishna et al.~2021}).

Five curricula had higher mean final evaluation loss than random initialization. Simpler tasks had a mean loss \(0.16\%\) lower than random initialization, while every synthetic curriculum had higher mean loss than the natural-text warmup. The evaluated reasoning and algorithmic probes showed no improvement from synthetic initialization.

\section{Protocol}\label{protocol}

The model used the Pythia-410M architecture and tokenizer with randomly initialized weights, a 12,288-token context, streamed public text datasets, and disjoint warmup, training, and evaluation slices (\href{https://arxiv.org/abs/2304.01373}{Biderman et al.~2023}). Every curriculum compared three conditions per seed: downstream training from random initialization, synthetic upstream training followed by the same downstream phase, and a natural-text upstream phase followed by that downstream phase. This common public-data protocol is a proxy comparison, not an exact reproduction of every source paper's corpus and training budget.

\publicationfigure[\linewidth]
  {images/pptrain-replication/claim_matrix.png}
  {/images/pptrain-replication/claim_matrix.png}
  {Claim matrix for the pptrain proxy run}
  {Claim matrix. A supported or contradicted comparison requires both the corresponding mean direction and agreement from at least two of three seeds; grey cells are inconclusive or unevaluated.}

\section{Results}\label{results}

\publicationfigure[\linewidth]
  {images/pptrain-replication/transfer_gap_vs_scratch.png}
  {/images/pptrain-replication/transfer_gap_vs_scratch.png}
  {Synthetic transfer gap versus random-init baseline}
  {Mean percentage reduction in final evaluation loss relative to random initialization. Positive values denote lower loss; only Simpler tasks has a positive mean.}

\publicationfigure[\linewidth]
  {images/pptrain-replication/compute_matched_nlp_prepretraining_gap.png}
  {/images/pptrain-replication/compute_matched_nlp_prepretraining_gap.png}
  {Synthetic transfer gap versus NLP prepretraining}
  {Mean percentage reduction in final evaluation loss relative to natural-text prepretraining. All six means are negative.}

Table 1. Mean final downstream losses and relative gaps across three seeds. Lower loss is better; positive gaps favor synthetic initialization.

\begin{longtable}[]{@{}
  >{\raggedright\arraybackslash}p{(\linewidth - 18\tabcolsep) * \real{0.0789}}
  >{\raggedright\arraybackslash}p{(\linewidth - 18\tabcolsep) * \real{0.0789}}
  >{\raggedleft\arraybackslash}p{(\linewidth - 18\tabcolsep) * \real{0.1053}}
  >{\raggedleft\arraybackslash}p{(\linewidth - 18\tabcolsep) * \real{0.1053}}
  >{\raggedleft\arraybackslash}p{(\linewidth - 18\tabcolsep) * \real{0.1053}}
  >{\raggedleft\arraybackslash}p{(\linewidth - 18\tabcolsep) * \real{0.1053}}
  >{\raggedleft\arraybackslash}p{(\linewidth - 18\tabcolsep) * \real{0.1053}}
  >{\raggedleft\arraybackslash}p{(\linewidth - 18\tabcolsep) * \real{0.1053}}
  >{\raggedleft\arraybackslash}p{(\linewidth - 18\tabcolsep) * \real{0.1053}}
  >{\raggedleft\arraybackslash}p{(\linewidth - 18\tabcolsep) * \real{0.1053}}@{}}
\toprule\noalign{}
\begin{minipage}[b]{\linewidth}\raggedright
Task
\end{minipage} & \begin{minipage}[b]{\linewidth}\raggedright
Synthetic preset
\end{minipage} & \begin{minipage}[b]{\linewidth}\raggedleft
Random-init loss
\end{minipage} & \begin{minipage}[b]{\linewidth}\raggedleft
Synthetic-init loss
\end{minipage} & \begin{minipage}[b]{\linewidth}\raggedleft
Text-warmup loss
\end{minipage} & \begin{minipage}[b]{\linewidth}\raggedleft
Synthetic vs random
\end{minipage} & \begin{minipage}[b]{\linewidth}\raggedleft
Synthetic vs text
\end{minipage} & \begin{minipage}[b]{\linewidth}\raggedleft
Random-init ppl
\end{minipage} & \begin{minipage}[b]{\linewidth}\raggedleft
Synthetic-init ppl
\end{minipage} & \begin{minipage}[b]{\linewidth}\raggedleft
Text-warmup ppl
\end{minipage} \\
\midrule\noalign{}
\endhead
\bottomrule\noalign{}
\endlastfoot
NCA & \texttt{paper\_web\_text} & 6.232 & 6.779 & 6.044 & -8.78\% & -12.17\% & 508.9 & 891.7 & 421.5 \\
LIME & \texttt{paper\_benchmark\_100k} & 5.134 & 6.032 & 4.713 & -17.60\% & -27.99\% & 170.8 & 419.3 & 111.4 \\
Simpler tasks & \texttt{paper\_unary\_core\_100k} & 6.237 & 6.227 & 6.040 & +0.16\% & -3.08\% & 511.1 & 506.1 & 420.1 \\
Procedural & \texttt{paper\_set\_len64} & 6.236 & 6.540 & 6.042 & -4.87\% & -8.24\% & 510.8 & 692.4 & 420.8 \\
Dyck & \texttt{paper\_k64} & 6.242 & 6.583 & 6.050 & -5.46\% & -8.79\% & 513.8 & 726.5 & 424.3 \\
Summarization & \texttt{paper\_ourtasks\_subset\_100k} & 5.828 & 6.623 & 5.559 & -13.66\% & -19.14\% & 339.6 & 766.1 & 259.7 \\
\end{longtable}

For NCA, held-out next-patch token accuracy was \(0.0024\%\); mean downstream cross-entropy was \(6.779\), compared with \(6.232\) from random initialization and \(6.044\) after natural-text warmup; and mean activation effective rank was \(4.6\), compared with \(48.4\) from random initialization. The generator passed the library's reference-parity fixtures, so the failure is not attributable to the generator outputs covered by those fixtures.

\publicationfigure[\linewidth]
  {images/pptrain-replication/nca_showcase.png}
  {/images/pptrain-replication/nca_showcase.png}
  {NCA upstream accuracy, downstream loss, and representation diagnostics}
  {NCA upstream accuracy, downstream loss, and representation measurements for the two controls and synthetic initialization.}

\section{Interpretation}\label{interpretation}

The random-initialization control asks whether the synthetic stage improves downstream training at all. The natural-text control asks whether the upstream updates are better spent on synthetic data or text. A loss of network plasticity during synthetic training, followed by distribution shift to natural text, could produce the observed ordering; related work documents cases of harmful warm starts, loss of plasticity, and feature distortion during shifted fine-tuning (\href{https://arxiv.org/abs/1910.08475}{Ash and Adams 2020}; \href{https://arxiv.org/abs/2306.13812}{Dohare et al.~2023}; \href{https://arxiv.org/abs/2202.10054}{Kumar et al.~2022}). This experiment does not distinguish that hypothesis from other causes.

Under this Pythia-410M public-data protocol, none of the six synthetic curricula produced lower mean final loss than natural-text prepretraining, and five produced higher mean loss than random initialization. Future transfer claims should report all three checks: comparison with random initialization, comparison with an equal natural-text upstream phase, and performance on the synthetic task itself.

\end{document}
